\documentclass[11pt]{article}
\usepackage[margin=1in]{geometry}
\usepackage{amsmath,amssymb,amsthm,mathtools}
\usepackage[hidelinks]{hyperref}

\newtheorem{theorem}{Theorem}[section]
\newtheorem{corollary}[theorem]{Corollary}
\newtheorem{lemma}[theorem]{Lemma}
\newtheorem{proposition}[theorem]{Proposition}
\theoremstyle{definition}
\newtheorem{definition}[theorem]{Definition}
\theoremstyle{remark}
\newtheorem{remark}[theorem]{Remark}

\newcommand{\Q}{Q}
\newcommand{\iotaR}{\iota_r}
\newcommand{\CAN}{\operatorname{CAN}}
\newcommand{\Kcov}{K}
\newcommand{\Vol}{V}

\title{Subcube isolation in hypercubes is radius-covering:\\
a coding-theoretic dictionary, perfect-code obstruction, and a phase transition}
\author{Working draft}
\date{September 2026}

\begin{document}
\maketitle

\begin{abstract}
Bre\v{s}ar and Rall asked whether the minimum number of vertices whose closed
neighborhoods destroy every copy of $Q_k$ in $Q_n$ always equals
$\gamma(Q_{n-k})$.  We identify the invariant exactly with a binary
radius-covering-array number.  This corrects the provenance of the smallest
counterexample, whose numerical value is established covering-array prior art,
and yields structural consequences not visible in the graph formulation.  We
prove a sharp fixed-codimension transition: the parameter is constant in
codimensions at most three and grows logarithmically from codimension four
onward.  We also derive a perfect-code projection obstruction.  For every
binary Hamming length $m=2^t-1$, $t\ge 3$, it implies
$\iota(Q_{m+k},Q_k)=\gamma(Q_m)$ exactly when $k=1$, giving infinitely many
strict inequalities while the surviving forbidden cube is fixed at $Q_2$.
The coordinate-copy theorem, the graph-to-array bridge, and the arithmetic core
are accompanied by Lean verification; the full packing obstruction remains an
explicit formalization target.
\end{abstract}

\section{Introduction}

Let $Q_n$ be the graph on $\{0,1\}^n$ with edges joining words at Hamming
distance one.  For graphs $G$ and $H$, the $H$-isolation number $\iota(G,H)$
is the minimum cardinality of a set $D\subseteq V(G)$ such that
$G-N[D]$ contains no copy of $H$.  Bre\v{s}ar and Rall asked whether
\begin{equation}\label{eq:BR}
  \iota(Q_n,Q_k)=\gamma(Q_{n-k})
\end{equation}
for every $0<k<n$.

The key observation is that the projected domination condition occurring in
this problem is not a new finite object: it is precisely a radius-covering
array.  The translation makes the proposed equality
\[
  \CAN_1(m,n,2)=\Kcov_2(m,1), \qquad m=n-k,
\]
where $\Kcov_q(m,r)$ is the minimum size of a $q$-ary radius-$r$ covering code.

\section{Coordinate copies of hypercubes}

\begin{lemma}[Coordinate-copy lemma]\label{lem:coordinate-copy}
Every subgraph of $Q_n$ isomorphic to $Q_k$ is a coordinate $k$-subcube.
Equivalently, after translating one vertex, its embedding is obtained by
injecting the $k$ intrinsic coordinate directions into $k$ distinct ambient
coordinates.
\end{lemma}

\begin{proof}
Color every ambient edge by the unique coordinate that it flips.  Every
four-cycle in a hypercube alternates between two distinct colors, so opposite
edges have the same color.  In $Q_k$, all edges in one intrinsic direction are
connected through a sequence of opposite-edge relations in four-cycles;
therefore they have one common ambient color.  Edges in distinct intrinsic
directions meet at a vertex and hence receive distinct colors.  Starting from
one image vertex, every domain word is therefore mapped by independently
flipping the corresponding fixed ambient coordinates.
\end{proof}

\section{The radius-covering dictionary}

For $q\ge2$, an $N\times n$ $q$-ary array has radius-covering strength $m$ and
radius $r$ if every choice of $m$ columns has covering radius at most $r$ in
$[q]^m$.  Write $\CAN_r(m,n,q)$ for the least possible number of rows.

Define $\iota_r(Q_n,Q_k)$ to be the minimum size of $D\subseteq Q_n$ such that
$Q_n-N_r[D]$ contains no copy of $Q_k$.  Thus $\iota_1=\iota$.

\begin{theorem}[Exact dictionary]\label{thm:dictionary}
For $0\le m\le n$ and $r\ge0$,
\[
  \iota_r(Q_n,Q_{n-m})=\CAN_r(m,n,2).
\]
In particular,
\[
  \iota(Q_n,Q_k)=\CAN_1(n-k,n,2).
\]
\end{theorem}

\begin{proof}
By Lemma~\ref{lem:coordinate-copy}, it suffices to hit every coordinate
$(n-m)$-face.  Such a face is specified by a set $S$ of $m$ fixed coordinates
and a pattern $a\in\{0,1\}^S$.  The minimum Hamming distance from a word $d$ to
that face is exactly $d_H(d|_S,a)$, because the free coordinates may be chosen
to agree with $d$.  Hence the radius-$r$ neighborhood of $D$ meets every such
face exactly when every $m$-column projection of the row array $D$ has covering
radius at most $r$.
\end{proof}

\section{A fixed-codimension phase transition}

Let
\[
  B_{m,r}=\sum_{j=0}^r {m\choose j}.
\]

\begin{theorem}[Phase transition]\label{thm:phase}
Fix $m\ge1$ and $r\ge0$.  As $n\to\infty$,
\[
\CAN_r(m,n,2)=
\begin{cases}
1, & m\le r,\\
2, & r<m\le 2r+1,\\
\Theta_{m,r}(\log n), & m\ge 2r+2.
\end{cases}
\]
More explicitly, in the third regime every $M$-row array satisfies
\[
  n\le (m-1)2^M,
\]
and a random construction gives
\[
 \CAN_r(m,n,2)
 \le 1+\left\lceil
 \frac{\log\!\left(2^m{n\choose m}\right)}
 {-\log\!\left(1-B_{m,r}/2^m\right)}
 \right\rceil.
\]
\end{theorem}

\begin{proof}
If $m\le r$, one row covers every $m$-tuple.  If $r<m\le2r+1$, one row is
insufficient, while the two constant complementary rows cover every target
because $\min\{|a|,m-|a|\}\le\lfloor m/2\rfloor\le r$.

Now assume $m\ge2r+2$.  Regard each column of an $M\times n$ array as a word in
$\{0,1\}^M$.  If one column type occurred at least $m$ times, select those
columns and a target with exactly $r+1$ ones.  Every projected row would be
$0^m$ or $1^m$, at distance at least $r+1$ from the target.  Thus every one of
the $2^M$ column types occurs at most $m-1$ times, proving the lower bound.

For the upper bound, choose $M$ independent uniform binary rows.  For a fixed
$m$-set of columns and a fixed target, the probability that no row lies within
radius $r$ is $(1-B_{m,r}/2^m)^M$.  A union bound over
$2^m{n\choose m}$ choices is less than one for the displayed value of $M$.
\end{proof}

\begin{corollary}
For fixed $m$,
\[
\iota(Q_n,Q_{n-m})=
\begin{cases}
1,&m=1,\\
2,&m=2,3,\\
\Theta_m(\log n),&m\ge4.
\end{cases}
\]
Consequently, for every fixed $m\ge4$,
$\iota(Q_n,Q_{n-m})/\gamma(Q_m)\to\infty$.
\end{corollary}

\section{Perfect projections and robust extension}

Let
\[
 \Vol_q(N,R)=\sum_{j=0}^R {N\choose j}(q-1)^j
\]
be the volume of a $q$-ary Hamming ball.

\begin{theorem}[Perfect-code extension obstruction]\label{thm:perfect}
Suppose a perfect $q$-ary radius-$r$ covering code of length $m$ exists, so
$\Kcov_q(m,r)\Vol_q(m,r)=q^m$.  Let $\ell\ge1$.  If
\[
 \Kcov_q(m,r)\Vol_q\!\left(m+\ell,
     r+\left\lfloor\ell/2\right\rfloor\right)>q^{m+\ell},
\]
then
\[
 \CAN_r(m,m+\ell,q)>\Kcov_q(m,r).
\]
\end{theorem}

\begin{proof}
Assume an array with $\Kcov_q(m,r)$ rows exists.  Every $m$-column projection
is a radius-$r$ covering code of optimal perfect size.  Since the sum of the
sizes of its radius-$r$ balls equals the whole space, those balls partition the
space; thus each projection has minimum distance at least $2r+1$.

Take two full rows at distance $d$.  Delete $\ell$ coordinates containing as
many of their disagreements as possible.  If $d\le2r+\ell$, the remaining
$m$-projection has distance at most $2r$, contradiction.  Therefore the full
array has minimum distance at least $2r+1+\ell$.  Hamming balls of radius
$r+\lfloor\ell/2\rfloor$ around its rows are disjoint, and the Hamming packing
bound contradicts the displayed strict inequality.
\end{proof}

\section{Hamming families}

\begin{corollary}[Binary fixed-$Q_2$ family]\label{cor:binary-hamming}
Let $t\ge3$ and $m=2^t-1$.  For every $k\ge1$,
\[
  \iota(Q_{m+k},Q_k)=\gamma(Q_m)
  \quad\Longleftrightarrow\quad k=1.
\]
In particular, $\iota(Q_{2^t+1},Q_2)>\gamma(Q_{2^t-1})$ for every $t\ge3$.
\end{corollary}

\begin{proof}
A binary Hamming code is a perfect radius-one code of length $m$.  Equality for
$k=1$ is the binary parity-extension theorem for radius-covering arrays.  For
$k=2$, Theorem~\ref{thm:perfect} applies because
\[
 \Vol_2(m+2,2)-4\Vol_2(m,1)=\frac{m(m-3)}2>0.
\]
Monotonicity in the number of columns gives strict inequality for every
$k\ge2$.
\end{proof}

\begin{corollary}[q-ary Hamming family]
Let $q$ be a prime power, $t\ge3$, and
$m=(q^t-1)/(q-1)$.  Then
\[
 \CAN_1(m,m+2,q)>\Kcov_q(m,1).
\]
\end{corollary}

\begin{proof}
The $q$-ary Hamming code is perfect, and
\[
 \Vol_q(m+2,2)-q^2\Vol_q(m,1)
 =\frac{(q-1)^2m(m+1-2q)}2>0.
\]
\end{proof}

\section{The smallest illustration and provenance}

The established radius-covering-array tables give
$\CAN_1(4,6,2)=5$.  Theorem~\ref{thm:dictionary} therefore yields
\[
  \iota(Q_6,Q_2)=5>4=\gamma(Q_4).
\]
One minimum set is
\[
 \{000000,000011,000101,111001,111110\}.
\]
This numerical value is not claimed as new.  An independent proof, exhaustive
Python enumeration, and a Lean finite certificate are retained for direct
auditing in graph language.

\section{Formal verification boundary}

The accompanying Lean development verifies the coordinate-copy theorem, the
exact distance-to-face identity, the equivalence between face hitting and the
radius-covering-array predicate, and the arithmetic core of the binary Hamming
obstruction.  The finite $(6,2)$ certificate is complete, but its exhaustive
four-row lower bound uses \texttt{native\_decide} with the generated axiom
reported explicitly.  A full formal proof of Theorem~\ref{thm:perfect},
including the general Hamming packing count, remains a stated gate rather than
a completed claim.

\section{Open directions}

The dictionary converts several graph questions into coding/design questions.
Natural next problems are to determine sharp logarithmic constants for fixed
$m,r$, characterize all pairs for which equality with a covering-code number
survives, and study analogous isolation parameters in Hamming graphs and other
Cartesian products.

\begin{thebibliography}{9}

\bibitem{BresarRall2026}
B.~Bre\v{s}ar and D.~F. Rall,
\newblock On the isolation numbers in graph products,
\newblock arXiv:2608.25752, 2026.

\bibitem{CKRSSP2010}
C.~J. Colbourn, G.~K\'eri, P.~P. Rivas Soriano, and J.-C. Schlage-Puchta,
\newblock Covering and radius-covering arrays: Constructions and classification,
\newblock \emph{Discrete Applied Mathematics} 158 (2010), 1158--1180.

\end{thebibliography}

\end{document}
